% !TeX encoding = UTF-8
% Reviewed conversion; see reports/revision-log.docx and bibliography.json.
\SourceChapter{5}{فصل پنجم: نتیجه‌گیری، جمع‌بندی و پیشنهادها}
% SOURCE Introduction_Chapter5.txt:1-1 [heading]
\SourceHeading{1}{۵{-}۱. مقدمه \lr{(Introduction)}}

% SOURCE Introduction_Chapter5.txt:2-2 [spacing]


% SOURCE Introduction_Chapter5.txt:3-3 [paragraph]
پژوهش حاضر با هدف طراحی، پیاده‌سازی و اعتبارسنجی یک چارچوب هیبریدی هوشمند و چندلایه‌ای جهت ارتقای امنیت، محرمانگی، دست‌نخورده ماندن بافت‌های تشخیصی و قابلیت انتقال زمان‌واقعی تصاویر پزشکی در شبکه‌های اینترنت اشیای پزشکی \lr{(IoMT)} و سامانه‌های سلامت الکترونیک \lr{(E{-}Health)} تدوین گردید. انتقال داده‌های حساس بالینی در بسترهای ارتباطی ناامن ناگزیر با تهدیدات سایبری متعددی از جمله دستکاری اطلاعات، حملات تحلیل آماری، حملات تفاضلی، شنود ناخواسته و افشای فراداده‌های هویتی بیماران مواجه است \ThesisCite{1}, \ThesisCite{2}, \ThesisCite{4}. چالش بنیادین در رمزشناسی تصویر پزشکی، ایجاد توازنی آرمانی میان سه ضلع یک سه‌گانه متعارض \lr{(Trilemma)} است: ۱) امنیت و پیچیدگی فوق‌العاده رمزنگاری جهت خنثی‌سازی کامل حملات رمزشناسی، ۲) عدم ایجاد کوچک‌ترین اعوجاج دیداری یا نویز در ناحیه حساس تشخیصی \lr{(ROI)} جهت صیانت مطلق از اصالت بالینی، و ۳) پایین بودن پیچیدگی محاسباتی جهت امکان پیاده‌سازی زمان‌واقعی بر روی سخت‌افزارهای کم‌مصرف پردازش لبه \lr{(Edge Devices)} \ThesisCite{1}, \ThesisCite{9}, \ThesisCite{10}.\par

% SOURCE Introduction_Chapter5.txt:4-4 [spacing]


% SOURCE Introduction_Chapter5.txt:5-5 [paragraph]
فرمول‌بندی ریاضی این سه‌گانه متعارض و تابع هدف بهینه‌سازی سیستم در رابطه ۵{-}۱ ارائه گردیده است \ThesisCite{1}, \ThesisCite{4}, \ThesisCite{8}.\par

% SOURCE Introduction_Chapter5.txt:6-6 [spacing]


% SOURCE Introduction_Chapter5.txt:7-7 [image]
\SourcePicture{eq5_0_trilemma_tradeoff.png}{رابطه ۵{-}۱: فرمول‌بندی ریاضی سه‌گانه متعارض رمزشناسی تصویر پزشکی و تابع هدف بهینه‌سازی}{[رابطه ۵{-}۱: فرمول‌بندی ریاضی سه‌گانه متعارض رمزشناسی تصویر پزشکی و تابع هدف بهینه‌سازی {-} \lr{eq5\_0\_trilemma\_tradeoff.png}]}{2}

% SOURCE Introduction_Chapter5.txt:8-8 [spacing]


% SOURCE Introduction_Chapter5.txt:9-9 [paragraph]
در این فصل، سنتز نهایی دستاوردهای تجربی، ارزیابی دقیق و اثبات قطعی فرضیات سه‌گانه پایان‌نامه بر پایه نتایج عددی فصل چهارم، تحلیل نحوه پوشش شکاف‌های پژوهشی استخراج‌شده در فصل دوم، تبیین محدودیت‌های واقعی فنی، و ارائه پیشنهادهای عملیاتی برای توسعه‌های آتی پژوهش ارائه می‌گردد \ThesisCite{1}, \ThesisCite{2}, \ThesisCite{3}, \ThesisCite{5}, \ThesisCite{9}.\par

% SOURCE Introduction_Chapter5.txt:10-10 [spacing]

% SOURCE Synthesis_Findings_Chapter5.txt:1-1 [heading]
\SourceHeading{1}{۵{-}۲. جمع‌بندی و سنتز دستاوردهای پژوهش \lr{(Synthesis of Research Findings)}}

% SOURCE Synthesis_Findings_Chapter5.txt:2-2 [spacing]


% SOURCE Synthesis_Findings_Chapter5.txt:3-3 [paragraph]
جهت دست‌یابی به اهداف تعیین‌شده، معماری امنیتی پیشنهادی در ۴ فاز یکپارچه طراحی و بر روی ۴ مجموعه‌داده مرجع بین‌المللی با مدالیته‌های متنوع تصویربرداری (شامل \lr{DRIVE}، \lr{RITE}، \lr{BraTS 2020}، و \lr{COVID{-}19 CXR}) پیاده‌سازی و ارزیابی شد \ThesisCite{1}, \ThesisCite{4}, \ThesisCite{9}. سنتز دستاوردهای کلیدی فازهای چهارگانه سیستم به شرح زیر است:\par

% SOURCE Synthesis_Findings_Chapter5.txt:4-4 [spacing]


% SOURCE Synthesis_Findings_Chapter5.txt:5-5 [paragraph]
۱. \textbf{فاز استخراج هوشمند \lr{ROI} و تولید کلید پویا \lr{(U{-}Net \& SHA{-}256 Key Generation)}:}\par

% SOURCE Synthesis_Findings_Chapter5.txt:6-6 [paragraph]
با به‌کارگیری شبکه عصبی عمیق \lr{U{-}Net}، ناحیه حساس تشخیصی \lr{(ROI)} به صورت خودکار قطعه‌بندی شد. سپس با استخراج گشتاورهای آماری این ناحیه و اعمال تابع هَش \lr{SHA{-}256}، کلید هَش ۲۵۶ بیتی $K_{\mathrm{hash}}$ به صورت کاملاً تصویر{-}پایه و پویا تولید گردید. این امر موجب شد که هر تصویر پزشکی کلید منحصر‌به‌فرد خود را داشته باشد و حملات متن‌واضح برگزیده \lr{(CPA)} و معلوم \lr{(KPA)} به صورت کامل خنثی شوند \ThesisCite{1}, \ThesisCite{2}.\par

% SOURCE Synthesis_Findings_Chapter5.txt:7-7 [spacing]


% SOURCE Synthesis_Findings_Chapter5.txt:8-8 [paragraph]
۲. \textbf{فاز تولید توالی‌های ابرآشوبی ۵بعدی \lr{(5D Hyperchaotic Sequence Generation)}:}\par

% SOURCE Synthesis_Findings_Chapter5.txt:9-9 [paragraph]
با استفاده از نگاشت ۵بعدی ابرآشوبی با دو توان مثبت لیپانوف ($L_1 = 1.24,\allowbreak{} L_2 = 0.85$)، توالی‌های آشوبی فوق‌العاده حساس به شرایط اولیه تولید گردیدند. کلید $K_{\mathrm{hash}}$ متغیرهای حالت اولیه سیستم را کنترل نموده و فضای کلید عظیم $2^{512}$ را ایجاد نمود که مقاومت مطلق در برابر حملات جستجوی فراگیر \lr{(Brute{-}Force)} را تضمین می‌کند \ThesisCite{1}, \ThesisCite{6}.\par

% SOURCE Synthesis_Findings_Chapter5.txt:10-10 [spacing]


% SOURCE Synthesis_Findings_Chapter5.txt:11-11 [paragraph]
۳. \textbf{فاز جایگشت زیگ‌زاگی و انتشار پویای \lr{DNA} \lr{(Zigzag Permutation \& Dynamic DNA Diffusion)}:}\par

% SOURCE Synthesis_Findings_Chapter5.txt:12-12 [paragraph]
با ترکیب جایگشت متقاطع زیگ‌زاگی \lr{(Confusion)} و قوانین انتشار مبتنی بر رمزنگاری \lr{DNA} \lr{(Diffusion)}، همبستگی خطی پیکسل‌های مجاور در ۳ جهت افقی، عمودی و قطری از $0.9838$ به کمتر از \textbf{$0.0008$} تنزل یافت و آنتروپی شانون به حد آرمانی \textbf{$7.9985$ بیت} رسید \ThesisCite{1}, \ThesisCite{4}.\par

% SOURCE Synthesis_Findings_Chapter5.txt:13-13 [spacing]


% SOURCE Synthesis_Findings_Chapter5.txt:14-14 [paragraph]
۴. \textbf{فاز پنهان‌نگاری آگاه به برجستگی بصری \lr{(Saliency{-}aware Steganography)}:}\par

% SOURCE Synthesis_Findings_Chapter5.txt:15-15 [paragraph]
فراداده هویتی بیمار پس از فشرده‌سازی بدون اتلاف \lr{zlib} (کاهش ۶۵\SourceGlyph{066A} حجم)، صرفاً در ۲ بیت کم‌ارزش \lr{(2{-}bit LSB)} پس‌زمینه کم‌برجستگی جای‌دهی شد. این استراتژی موجب شد که ناحیه تشخیصی \lr{ROI} به صورت ۱۰۰\SourceGlyph{066A} دست‌نخورده باقی بماند ($\mathrm{PSNR}_{\mathrm{ROI}} = \infty$) و شفافیت استگو کل تصویر به \textbf{$\mathrm{PSNR} = 52.41\mathrm{\ dB}$} و \textbf{$\mathrm{SSIM} = 0.9982$} برسد \ThesisCite{3}, \ThesisCite{8}, \ThesisCite{9}.\par

% SOURCE Synthesis_Findings_Chapter5.txt:16-16 [spacing]


% SOURCE Synthesis_Findings_Chapter5.txt:17-17 [paragraph]
فرمول‌بندی ریاضی و سنتز روابط چهارگانه فازهای سیستم در رابطه ۵{-}۲ ارائه گردیده است \ThesisCite{1}, \ThesisCite{4}, \ThesisCite{6}, \ThesisCite{8}.\par

% SOURCE Synthesis_Findings_Chapter5.txt:18-18 [spacing]


% SOURCE Synthesis_Findings_Chapter5.txt:19-19 [image]
\SourcePicture{eq5_2_synthesis_4phases_math.png}{رابطه ۵{-}۲: فرمول‌بندی ریاضی و سنتز فازهای چهارگانه معماری هیبریدی}{[رابطه ۵{-}۲: فرمول‌بندی ریاضی و سنتز فازهای چهارگانه معماری هیبریدی {-} \lr{eq5\_2\_synthesis\_4phases\_math.png}]}{2}

% SOURCE Synthesis_Findings_Chapter5.txt:20-20 [spacing]

% SOURCE Hypotheses_Verification_Chapter5.txt:1-1 [heading]
\SourceHeading{1}{۵{-}۳. ارزیابی و اثبات تجربی فرضیات سه‌گانه پایان‌نامه \lr{(Empirical Verification of Thesis Hypotheses)}}

% SOURCE Hypotheses_Verification_Chapter5.txt:2-2 [spacing]


% SOURCE Hypotheses_Verification_Chapter5.txt:3-3 [paragraph]
ارزیابی کمّی، سنجش ریاضی و تحلیل‌های آماری صورت‌گرفته در فصل چهارم بر روی ۴ مجموعه‌داده مرجع بین‌المللی \lr{(DRIVE, RITE, BraTS 2020, COVID{-}19 CXR)} و بستر پردازش لبه \lr{NVIDIA Jetson TX2}، صحت کامل، اعتبار روش‌شناختی و تحقق ۱۰۰ درصدی فرضیات سه‌گانه تدوین‌شده در فصل اول (بخش ۱{-}۴) را به صورت تجربی و ریاضی به اثبات رساند \ThesisCite{1}, \ThesisCite{2}, \ThesisCite{4}, \ThesisCite{9}, \ThesisCite{11}.\par

% SOURCE Hypotheses_Verification_Chapter5.txt:4-4 [spacing]


% SOURCE Hypotheses_Verification_Chapter5.txt:5-5 [paragraph]
جدول ۵{-}۱ ماتریس جامع اثبات تجربی فرضیات سه‌گانه پایان‌نامه را بر پایه مقایسه شاخص‌های کمّی هدف‌گذاری‌شده در فصل اول و دستاوردهای عددی واقعی حاصل در فصل چهارم نشان می‌دهد.\par

% SOURCE Hypotheses_Verification_Chapter5.txt:6-6 [spacing]


% SOURCE Hypotheses_Verification_Chapter5.txt:7-19 [table]
\par\addvspace{\baselineskip}
\begingroup\fontsize{12}{15}\selectfont\setlength{\tabcolsep}{4pt}
\renewcommand{\arraystretch}{1.25}
\SourceCaption{lot}{جدول ۵{-}۱: ماتریس جامع اثبات تجربی فرضیات سه‌گانه پایان‌نامه بر پایه دستاوردهای کمّی فصل چهارم}
\begin{NoHyper}\begin{longtable}{@{}>{\raggedleft\arraybackslash}p{\dimexpr(\linewidth-8\tabcolsep)/4\relax}>{\raggedleft\arraybackslash}p{\dimexpr(\linewidth-8\tabcolsep)/4\relax}>{\raggedleft\arraybackslash}p{\dimexpr(\linewidth-8\tabcolsep)/4\relax}>{\raggedleft\arraybackslash}p{\dimexpr(\linewidth-8\tabcolsep)/4\relax}@{}}
\toprule
فرضیه پژوهش (مندرج در فصل اول) & شاخص‌های کمّی هدف‌گذاری‌شده & مقادیر عددی تجربی حاصل در فصل چهارم & وضعیت اثبات فرضیه \\
\midrule\endfirsthead
فرضیه پژوهش (مندرج در فصل اول) & شاخص‌های کمّی هدف‌گذاری‌شده & مقادیر عددی تجربی حاصل در فصل چهارم & وضعیت اثبات فرضیه \\\midrule\endhead
\textbf{فرضیه اول ($H_1$):} پویایی کلید و مقاومت تفاضلی & $\mathrm{NPCR} > 99.5693\%$ & $\mathrm{NPCR} = 99.6254\%$ & \textbf{تایید قطعی \lr{(PASSED)}} \\
 & $\mathrm{UACI} \in [33.31\%,\allowbreak{} 33.61\%]$ & $\mathrm{UACI} = 33.5120\%$ &  \\
\textbf{فرضیه دوم ($H_2$):} تصادفی‌سازی و آنتروپی آرمانی & $H(m) \approx 8.00\text{ bits}$ & $H(m) = 7.9985\text{ bits}$ & \textbf{تایید قطعی \lr{(PASSED)}} \\
 & $r_{xy} \approx 0.0000$ & $r_{xy} < 0.0008$ &  \\
 & $\text{NIST SP800-22 } p \ge 0.01$ & پاس کردن ۱۵ آزمون ($p \ge 0.0125$) &  \\
\textbf{فرضیه سوم ($H_3$):} شفافیت استگو و صیانت از \lr{ROI} & $\mathrm{PSNR}_{\mathrm{Total}} > 50\text{ dB}$ & $\mathrm{PSNR}_{\mathrm{Total}} = 52.41\text{ dB}$ & \textbf{تایید قطعی \lr{(PASSED)}} \\
 & $\mathrm{PSNR}_{\mathrm{ROI}} = \infty$ & $\mathrm{PSNR}_{\mathrm{ROI}} = \infty$ (عدم نویز) &  \\
 & $\mathrm{SSIM} \ge 0.9980$ & $\mathrm{SSIM} = 0.9982$ &  \\
\bottomrule
\end{longtable}\end{NoHyper}\endgroup
\par\addvspace{\baselineskip}

% SOURCE Hypotheses_Verification_Chapter5.txt:20-20 [spacing]


% SOURCE Hypotheses_Verification_Chapter5.txt:21-21 [paragraph]
تحلیل تفصیلی و سنتز اثبات تجربی هر یک از فرضیات سه‌گانه پایان‌نامه به شرح زیر است:\par

% SOURCE Hypotheses_Verification_Chapter5.txt:22-22 [spacing]


% SOURCE Hypotheses_Verification_Chapter5.txt:23-23 [paragraph]
۱. \textbf{ارزیابی و اثبات فرضیه اول ($H_1$ {-} پویایی کلید و مقاومت در برابر حملات تفاضلی و رمزشناسی):}\par

% SOURCE Hypotheses_Verification_Chapter5.txt:24-24 [paragraph]
مطابق فرضیه اول تدوین‌شده در بخش ۱{-}۴، ادعا گردید که استخراج ویژگی‌های آماری ناحیه تشخیصی \lr{ROI} و تولید کلید هَش پویا، موجب ایجاد حساسیت شدیدی نسبت به کوچک‌ترین تغییر در تصویر ورودی گشته و مقاومت تفاضلی سیستم را تضمین می‌نماید. نتایج تجربی فصل چهارم (بخش ۴{-}۶) اثبات نمود که با تغییر تنها ۱ پیکسل در تصویر ورودی، کلید هَش ۲۵۶ بیتی $K_{\mathrm{hash}}$ به صورت کامل دگرگون شده و متغیرهای حالت اولیه سیستم ۵بعدی ابرآشوبی را کاملاً تغییر می‌دهد \ThesisCite{1}, \ThesisCite{2}, \ThesisCite{6}. شبیه‌سازی‌های انجام‌شده بر روی ۱۰۰ جفت تصویر متناظر نشان داد که متوسط شاخص \lr{NPCR} به \textbf{۹۹.۶۲۵۴\SourceGlyph{066A}} (فراتر از حد آستانه تئوریک ۹۹.۵۶۹۳\SourceGlyph{066A}) و متوسط شاخص \lr{UACI} به \textbf{۳۳.۵۱۲۰\SourceGlyph{066A}} (در بازه آرمانی $[33.3115\%,\allowbreak{} 33.6156\%]$) رسیده و نرخ قبولی آزمون‌های تفاضلی بر روی تمامی ۱۰۰ جفت تصویر برابر با \textbf{۱۰۰\SourceGlyph{066A}} گردیده است. این نتایج، خنثی‌سازی کامل حملات متن‌واضح برگزیده \lr{(CPA)}، متن‌واضح معلوم \lr{(KPA)} و حملات تفاضلی را اثبات نموده و فرضیه اول را به صورت قطعی تایید می‌نماید \ThesisCite{1}, \ThesisCite{4}.\par

% SOURCE Hypotheses_Verification_Chapter5.txt:25-25 [spacing]


% SOURCE Hypotheses_Verification_Chapter5.txt:26-26 [paragraph]
۲. \textbf{ارزیابی و اثبات فرضیه دوم ($H_2$ {-} رفتار شبه‌تصادفی کامل، آنتروپی آرمانی و آزمون‌های \lr{NIST SP800{-}22}):}\par

% SOURCE Hypotheses_Verification_Chapter5.txt:27-27 [paragraph]
مطابق فرضیه دوم، ادعا شد که ترکیب جایگشت زیگ‌زاگی و انتشار پویای \lr{DNA} مبتنی بر نگاشت ۵بعدی ابرآشوبی، خروجی‌های رمزشده‌ای با رفتار شبه‌تصادفی کامل و آنتروپی نزدیک به ۸ بیت تولید خواهد نمود. نتایج تجربی فصل چهارم (بخش‌های ۴{-}۴، ۴{-}۵ و ۴{-}۹) نشان داد که متوسط آنتروپی اطلاعات شانون برابر با \textbf{۷.۹۹۸۵ بیت} (در مقایسه با حد نظری آرمانی ۸.۰۰ بیت) بوده و ضریب همبستگی پیکسل‌های مجاور در جهت‌های افقی، عمودی و قطری از $0.9838$ به کمتر از \textbf{$0.0008$} تنزل یافته است \ThesisCite{1}, \ThesisCite{5}, \ThesisCite{6}. همچنین، آزمون مجذور کای ($\chi^2 < 293.2478$) یکنواختی کامل هیستوگرام تصاویر رمزشده را تایید نمود. اعتبارسنجی نهایی با اجرای مجموعه آزمون‌های ۱۵گانه استاندارد \lr{NIST SP800{-}22} بر روی ۱۰۰ توالی $10^6$ بیتی صورت پذیرفت که تمامی $p\text{-value}$ها در بازه $[0.0125,\allowbreak{} 0.9850]$ قرار گرفته (فراتر از آستانه معناداری $\alpha = 0.01$) و نرخ قبولی توالی‌ها بین ۹۸\SourceGlyph{066A} تا ۱۰۰\SourceGlyph{066A} (بالاتر از آستانه $96.0\%$) ثبت گردید. این نتایج اثبات می‌نمایند که فرضیه دوم پژوهش با بالاترین سطح اطمینان آماری محقق شده است \ThesisCite{5}, \ThesisCite{8}, \ThesisCite{11}.\par

% SOURCE Hypotheses_Verification_Chapter5.txt:28-28 [spacing]


% SOURCE Hypotheses_Verification_Chapter5.txt:29-29 [paragraph]
۳. \textbf{ارزیابی و اثبات فرضیه سوم ($H_3$ {-} شفافیت دیداری استگو، صیانت ۱۰۰\SourceGlyph{066A} از بافت تشخیصی \lr{ROI} و پنهان‌نگاری فراداده):}\par

% SOURCE Hypotheses_Verification_Chapter5.txt:30-30 [paragraph]
مطابق فرضیه سوم، ادعا شد که ماژول پنهان‌نگاری آگاه به برجستگی بصری \lr{(Saliency{-}aware)} فراداده بیمار را بدون ایجاد کوچک‌ترین نویز در ناحیه \lr{ROI} جای‌دهی کرده و شفافیت استگو کل تصویر را با $\mathrm{PSNR} > 50\mathrm{\ dB}$ حفظ خواهد نمود. نتایج تجربی فصل چهارم (بخش ۴{-}۷) اثبات کرد که جای‌دهی ۲ بیتی \lr{LSB} فراداده فشرده‌شده \lr{zlib} صرفاً در پس‌زمینه کم‌برجستگی، موجب صیانت ۱۰۰ درصدی از بافت تشخیصی گردیده است (\textbf{$\mathrm{PSNR}_{\mathrm{ROI}} = \infty$} و $\mathrm{MSE}_{\mathrm{ROI}} = 0$). همچنین، متوسط شاخص کیفیت دیداری کل تصویر به \textbf{$\mathrm{PSNR} = 52.41\mathrm{\ dB}$}، شاخص شباهت ساختاری به \textbf{$\mathrm{SSIM} = 0.9982$} و ظرفیت جای‌دهی به \textbf{$1.18\mathrm{\ bpp}$} رسید که عدم امکان تشخیص دیداری تغییرات توسط متخصصان رادیولوژی را تضمین می‌کند \ThesisCite{3}, \ThesisCite{8}, \ThesisCite{9}. آزمایش‌های پایداری (بخش ۴{-}۸) نیز اثبات کردند که سیستم دارای ویژگی بازسازی‌پذیری خطاپذیر \lr{(Fault{-}Tolerant)} بوده و حتی تحت ۵۰\SourceGlyph{066A} قطع داده، بافت تشخیصی \lr{ROI} با کیفیت بالینی مناسب ($\mathrm{PSNR} > 32.0\mathrm{\ dB}$) بازیابی می‌شود. این دستاوردها فرضیه سوم پژوهش را نیز به صورت قطعی اثبات می‌نمایند \ThesisCite{5}, \ThesisCite{6}.\par

% SOURCE Hypotheses_Verification_Chapter5.txt:31-31 [spacing]


% SOURCE Hypotheses_Verification_Chapter5.txt:32-32 [paragraph]
فرمول‌بندی ریاضی ماتریس ارزیابی و اثبات تجربی فرضیات سه‌گانه در رابطه ۵{-}۳ نشان داده شده است.\par

% SOURCE Hypotheses_Verification_Chapter5.txt:33-33 [spacing]


% SOURCE Hypotheses_Verification_Chapter5.txt:34-34 [image]
\SourcePicture{eq5_1_hypotheses_matrix.png}{رابطه ۵{-}۳: فرمول‌بندی ریاضی ماتریس ارزیابی و اثبات تجربی فرضیات سه‌گانه پایان‌نامه}{[رابطه ۵{-}۳: فرمول‌بندی ریاضی ماتریس ارزیابی و اثبات تجربی فرضیات سه‌گانه پایان‌نامه {-} \lr{eq5\_1\_hypotheses\_matrix.png}]}{2}

% SOURCE Hypotheses_Verification_Chapter5.txt:35-35 [spacing]

% SOURCE Addressing_Gaps_Chapter5.txt:1-1 [heading]
\SourceHeading{1}{۵{-}۴. تحلیل پاسخگویی به پرسش‌های پژوهش و پوشش شکاف‌های علمی \lr{(Answering Research Questions \& Addressing Gaps)}}

% SOURCE Addressing_Gaps_Chapter5.txt:2-2 [spacing]


% SOURCE Addressing_Gaps_Chapter5.txt:3-3 [paragraph]
ارزیابی انتقادی و تحلیل سنتزیافته‌ی نتایج تجربی فصل چهارم نشان می‌دهد که معماری هیبریدی پیشنهادی به صورت دقیق و ساختارمند به تمامی پرسش‌های اصلی پژوهش (تدوین‌شده در بخش ۱{-}۴{-}۱) پاسخ داده و سه شکاف پژوهشی اساسی استخراج‌شده در ادبیات تحقیق (بخش ۲{-}۴) را برطرف ساخته است. تحلیل تفصیلی نحوه پوشش این شکاف‌های علمی بر پایه دستاوردهای عددی تجربی به شرح زیر ارائه می‌گردد \ThesisCite{1}, \ThesisCite{2}, \ThesisCite{4}, \ThesisCite{9}, \ThesisCite{10}, \ThesisCite{11}:\par

% SOURCE Addressing_Gaps_Chapter5.txt:4-4 [spacing]


% SOURCE Addressing_Gaps_Chapter5.txt:5-5 [paragraph]
۱. \textbf{پاسخ به پرسش و رفع شکاف اول (تولید کلید پویا و خنثی‌سازی کامل حملات متن‌واضح \lr{CPA/KPA}):}\par

% SOURCE Addressing_Gaps_Chapter5.txt:6-6 [paragraph]
{-} \textbf{پرسش علمی و شکاف موجود:} اکثر الگوریتم‌های رمزشناسی پیشین از کلیدهای ایستا یا وابسته به پارامترهای ثابت استفاده می‌نمودند که سیستم را در برابر حملات متن‌واضح برگزیده \lr{(CPA)} و متن‌واضح معلوم \lr{(KPA)} بسیار آسیب‌پذیر می‌ساخت \ThesisCite{2}, \ThesisCite{6}. چالش اصلی، چگونگی تولید کلید هَش پویا و وابسته به تصویر بدون ایجاد حساسیت مفرط به نویزهای پس‌زمینه بود.\par

% SOURCE Addressing_Gaps_Chapter5.txt:7-7 [paragraph]
{-} \textbf{راهکار و دستاورد تجربی این پژوهش:} با به‌کارگیری شبکه عصبی عمیق \lr{U{-}Net} جهت قطعه‌بندی خودکار ناحیه حساس تشخیصی \lr{(ROI)} و استخراج گشتاورهای آماری این ناحیه، تابع هَش \lr{SHA{-}256} کلید ۲۵۶ بیتی $K_{\mathrm{hash}}$ را به صورت کاملاً پویا و تصویر{-}پایه تولید نمود. با تغییر تنها ۱ پیکسل در تصویر ورودی، کلید $K_{\mathrm{hash}}$ به صورت کامل دگرگون شده و مقادیر شاخص‌های تفاضلی به $\mathrm{NPCR} = 99.6254\%$ و $\mathrm{UACI} = 33.5120\%$ رسید که آسیب‌پذیری در برابر حملات \lr{CPA} و \lr{KPA} را به صورت ۱۰۰\SourceGlyph{066A} خنثی ساخت \ThesisCite{1}, \ThesisCite{4}.\par

% SOURCE Addressing_Gaps_Chapter5.txt:8-8 [spacing]


% SOURCE Addressing_Gaps_Chapter5.txt:9-9 [paragraph]
۲. \textbf{پاسخ به پرسش و رفع شکاف دوم (رفع چالش تقابل ظرفیت پنهان‌نگاری و کیفیت تشخیصی بالینی {-} \lr{PSNR vs ROI}):}\par

% SOURCE Addressing_Gaps_Chapter5.txt:10-10 [paragraph]
{-} \textbf{پرسش علمی و شکاف موجود:} روش‌های پنهان‌نگاری پیشین، فراداده بیمار را بدون توجه به ارزش بالینی نواحی تصویر جای‌دهی می‌کردند که منجر به ایجاد نویز دیداری و اعوجاج آناتومیک در ناحیه \lr{ROI} می‌گردید ($\mathrm{PSNR} < 40\mathrm{\ dB}$) \ThesisCite{3}, \ThesisCite{9}, \ThesisCite{10}.\par

% SOURCE Addressing_Gaps_Chapter5.txt:11-11 [paragraph]
{-} \textbf{راهکار و دستاورد تجربی این پژوهش:} با ترکیب فشرده‌سازی بدون اتلاف \lr{zlib} (کاهش ۶۵ درصدی حجم فراداده) و الگوریتم تشخیص برجستگی بصری \lr{(Spectral Residual Saliency)}، فراداده هویتی بیمار صرفاً در ۲ بیت کم‌ارزش \lr{(2{-}bit LSB)} پس‌زمینه کم‌برجستگی \lr{(Non{-}ROI)} جای‌دهی شد. این استراتژی موجب گردید که ناحیه تشخیصی پزشکی دچار صفر درصد اعوجاج گردد ($\mathrm{PSNR}_{\mathrm{ROI}} = \infty$ و $\mathrm{MSE}_{\mathrm{ROI}} = 0$) و شفافیت استگو کل تصویر به شاخص بی‌نظیر $\mathrm{PSNR} = 52.41\mathrm{\ dB}$ و $\mathrm{SSIM} = 0.9982$ دست یابد \ThesisCite{3}, \ThesisCite{8}, \ThesisCite{9}.\par

% SOURCE Addressing_Gaps_Chapter5.txt:12-12 [spacing]


% SOURCE Addressing_Gaps_Chapter5.txt:13-13 [paragraph]
۳. \textbf{پاسخ به پرسش و رفع شکاف سوم (یکپارچه‌سازی لایه‌های امنیتی برای سخت‌افزارهای لبه \lr{IoMT} و کارایی زمان‌واقعی):}\par

% SOURCE Addressing_Gaps_Chapter5.txt:14-14 [paragraph]
{-} \textbf{پرسش علمی و شکاف موجود:} کارهای پیشین به دلیل پیچیدگی محاسباتی بالا، قادر به مجتمع‌سازی قطعه‌بندی عمیق، رمزنگاری ۵بعدی ابرآشوبی و پنهان‌نگاری فراداده در قالب یک سیستم سبک برای تجهیزات لبه \lr{IoMT} نبودند \ThesisCite{1}, \ThesisCite{2}, \ThesisCite{10}.\par

% SOURCE Addressing_Gaps_Chapter5.txt:15-15 [paragraph]
{-} \textbf{راهکار و دستاورد تجربی این پژوهش:} با بهینه‌سازی محاسباتی کدها و ارزیابی واقعی بر روی پلتفرم کم‌مصرف پردازش لبه \lr{NVIDIA Jetson TX2}، زمان اجرای کل سیستم برای یک تصویر پزشکی استاندارد به \textbf{۲.۸۵ ثانیه} رسید (تفکیک زمان اجرا: قطعه‌بندی \lr{U{-}Net} و کلید ۰.۸۲ ثانیه، سیستم ۵بعدی ۰.۴۵ ثانیه، زیگ‌زاگ و \lr{DNA} برابر ۰.۹۸ ثانیه، و \lr{zlib/Saliency} برابر ۰.۶۰ ثانیه). این زمان اجرا قابلیت استقرار زمان‌واقعی سیستم را در شبکه \lr{IoMT} و سلامت همراه به اثبات رساند \ThesisCite{10}, \ThesisCite{11}.\par

% SOURCE Addressing_Gaps_Chapter5.txt:16-16 [spacing]


% SOURCE Addressing_Gaps_Chapter5.txt:17-17 [paragraph]
فرمول‌بندی ریاضی پاسخگویی به پرسش‌های پژوهش و انطباق دستاوردهای تجربی با شکاف‌های علمی در رابطه ۵{-}۴ ارائه گردیده است \ThesisCite{1}, \ThesisCite{4}, \ThesisCite{9}, \ThesisCite{10}, \ThesisCite{11}.\par

% SOURCE Addressing_Gaps_Chapter5.txt:18-18 [spacing]


% SOURCE Addressing_Gaps_Chapter5.txt:19-19 [image]
\SourcePicture{eq5_3_gaps_mapping_math.png}{رابطه ۵{-}۴: فرمول‌بندی ریاضی پاسخگویی به پرسش‌های پژوهش و انطباق دستاوردهای تجربی با شکاف‌های علمی}{[رابطه ۵{-}۴: فرمول‌بندی ریاضی پاسخگویی به پرسش‌های پژوهش و انطباق دستاوردهای تجربی با شکاف‌های علمی {-} \lr{eq5\_3\_gaps\_mapping\_math.png}]}{2}
% SOURCE Methodological_Contributions_Chapter5.txt:1-1 [heading]
\SourceHeading{1}{۵{-}۵. دستاوردهای روش‌شناختی، نوآوری‌ها و کاربردهای بالینی در شبکه \lr{IoMT} \lr{(Methodological Contributions \& Clinical Applications)}}

% SOURCE Methodological_Contributions_Chapter5.txt:2-2 [spacing]


% SOURCE Methodological_Contributions_Chapter5.txt:3-3 [paragraph]
دستاوردهای روش‌شناختی، نوآوری‌های علمی و کاربردهای بالینی این پژوهش در جهت ارتقای امنیت انتقال داده‌های حساس در شبکه اینترنت اشیای پزشکی \lr{(IoMT)}، سامانه‌های سلامت همراه \lr{(mHealth)} و بیمارستان‌های هوشمند تدوین گردیده است \ThesisCite{1}, \ThesisCite{2}, \ThesisCite{10}. رمزشناسی تصویر پزشکی به دلیل ماهیت ویژه داده‌های بالینی، همواره با چالش صیانت مطلق از بافت‌های تشخیصی و محدودیت‌های سخت‌افزاری پردازش لبه مواجه بوده است \ThesisCite{3}, \ThesisCite{9}, \ThesisCite{11}. نوآوری‌ها و دستاوردهای چهارگانه اصلی این پژوهش به شرح زیر طبقه‌بندی می‌گردند:\par

% SOURCE Methodological_Contributions_Chapter5.txt:4-4 [spacing]


% SOURCE Methodological_Contributions_Chapter5.txt:5-5 [paragraph]
۱. \textbf{معماری یکپارچه چهارگانه و تلفیق هوشمند یادگیری عمیق با رمزشناسی آشوبی:}\par

% SOURCE Methodological_Contributions_Chapter5.txt:6-6 [paragraph]
تلفیق هوشمند قطعه‌بندی عمیق \lr{U{-}Net}، کلید پویا \lr{SHA{-}256}، سیستم ۵بعدی ابرآشوبی، انتشار پویای \lr{DNA}، و پنهان‌نگاری آگاه به برجستگی بصری \lr{(Saliency{-}aware)} برای نخستین بار در ادبیات تحقیق ارائه گردید \ThesisCite{1}, \ThesisCite{5}, \ThesisCite{8}. این معماری هیبریدی توازنی کامل میان سه ضلع متعارض امنیت، کیفیت دیداری و سرعت پردازش ایجاد نموده است \ThesisCite{1}, \ThesisCite{4}, \ThesisCite{6}.\par

% SOURCE Methodological_Contributions_Chapter5.txt:7-7 [spacing]


% SOURCE Methodological_Contributions_Chapter5.txt:8-8 [paragraph]
۲. \textbf{تولید کلید پویا بر پایه گشتاورهای آماری \lr{ROI} و خنثی‌سازی کامل حملات \lr{CPA/KPA}:}\par

% SOURCE Methodological_Contributions_Chapter5.txt:9-9 [paragraph]
بر خلاف سیستم‌های سنتی آشوبی که از کلیدهای ایستا استفاده می‌نمایند، این پژوهش کلید هَش ۲۵۶ بیتی $K_{\mathrm{hash}}$ را به صورت کاملاً تصویر{-}پایه از گشتاورهای ناحیه حساس تشخیصی استخراج نمود \ThesisCite{1}, \ThesisCite{2}. این نوآوری موجب ایجاد اثر بهمنی کامل، حساسیت تک‌پیکسلی، و شکست تمامی حملات متن‌واضح برگزیده \lr{(CPA)} و معلوم \lr{(KPA)} گردید \ThesisCite{1}, \ThesisCite{4}.\par

% SOURCE Methodological_Contributions_Chapter5.txt:10-10 [spacing]


% SOURCE Methodological_Contributions_Chapter5.txt:11-11 [paragraph]
۳. \textbf{پنهان‌نگاری آگاه به برجستگی بصری و صیانت ۱۰۰ درصدی از ناحیه \lr{ROI} ($\mathrm{PSNR}_{\mathrm{ROI}} = \infty$):}\par

% SOURCE Methodological_Contributions_Chapter5.txt:12-12 [paragraph]
با جای‌دهی فراداده فشرده‌شده \lr{zlib} بیمار صرفاً در پس‌زمینه کم‌برجستگی، ناحیه حساس تشخیصی به صورت ۱۰۰\SourceGlyph{066A} دست‌نخورده و بدون ایجاد کوچک‌ترین اعوجاج آناتومیک باقی ماند ($\mathrm{PSNR}_{\mathrm{ROI}} = \infty$ و $\mathrm{MSE}_{\mathrm{ROI}} = 0$) \ThesisCite{3}, \ThesisCite{8}, \ThesisCite{9}. این دستاورد امکان انتقال هم‌زمان تصویر تشخیصی و فراداده پرونده الکترونیک سلامت \lr{(EHR)} را در یک بستر یکپارچه فراهم می‌سازد \ThesisCite{2}, \ThesisCite{9}.\par

% SOURCE Methodological_Contributions_Chapter5.txt:13-13 [spacing]


% SOURCE Methodological_Contributions_Chapter5.txt:14-14 [paragraph]
۴. \textbf{قابلیت بازسازی‌پذیری خطاپذیر \lr{(Fault{-}Tolerant Reconstruction)} و استقرار زمان‌واقعی لبه:}\par

% SOURCE Methodological_Contributions_Chapter5.txt:15-15 [paragraph]
سیستم پیشنهادی امکان بازیابی بافت تشخیصی با کیفیت بالینی مناسب ($\mathrm{PSNR} > 32.0\mathrm{\ dB}$ و $\mathrm{SSIM}_{\mathrm{ROI}} \ge 0.9815$) را حتی پس از اعمال نویز شدید فلفل‌نمکی، گوسی و قطع ۵۰ درصدی داده‌ها \lr{(Cropping Attack)} در کانال ارتباطی فراهم می‌کند \ThesisCite{5}, \ThesisCite{6}. علاوه بر این، زمان اجرای کل سیستم بر روی پلتفرم کم‌مصرف پردازش لبه \lr{NVIDIA Jetson TX2} برابر با \textbf{۲.۸۵ ثانیه} ثبت شد که استقرار عملیاتی آن را در تجهیزات پزشکی \lr{IoMT} اثبات می‌نماید \ThesisCite{10}, \ThesisCite{11}.\par

% SOURCE Methodological_Contributions_Chapter5.txt:16-16 [spacing]


% SOURCE Methodological_Contributions_Chapter5.txt:17-17 [paragraph]
فرمول‌بندی ریاضی و سنتز دستاوردهای روش‌شناختی و کاربردهای بالینی در شبکه \lr{IoMT} در رابطه ۵{-}۵ ارائه گردیده است \ThesisCite{1}, \ThesisCite{5}, \ThesisCite{9}, \ThesisCite{10}, \ThesisCite{11}.\par

% SOURCE Methodological_Contributions_Chapter5.txt:18-18 [spacing]


% SOURCE Methodological_Contributions_Chapter5.txt:19-19 [image]
\SourcePicture{eq5_4_clinical_contributions_math.png}{رابطه ۵{-}۵: فرمول‌بندی ریاضی دستاوردهای روش‌شناختی و کاربردهای بالینی در شبکه \lr{IoMT}}{[رابطه ۵{-}۵: فرمول‌بندی ریاضی دستاوردهای روش‌شناختی و کاربردهای بالینی در شبکه \lr{IoMT {-} eq5\_4\_clinical\_contributions\_math.png}]}{2}

% SOURCE Methodological_Contributions_Chapter5.txt:20-20 [spacing]

% SOURCE Technical_Boundaries_Chapter5.txt:1-1 [heading]
\SourceHeading{1}{۵{-}۶. محدودیت‌های واقعی و فنی پژوهش \lr{(Actual Technical \& Practical Boundaries)}}

% SOURCE Technical_Boundaries_Chapter5.txt:2-2 [spacing]


% SOURCE Technical_Boundaries_Chapter5.txt:3-3 [paragraph]
علیرغم دست‌یابی به تمامی اهداف کمّی تعیین‌شده در طرح پژوهش و اثبات قطعی فرضیات سه‌گانه پایان‌نامه، شفافیت علمی و زبان احتیاط پژوهشی \lr{(Scientific Hedging)} ایجاب می‌نماید که مرزها، چالش‌ها و محدودیت‌های واقعی فنی الگوریتم پیشنهادی به روشنی تبیین گردند \ThesisCite{1}, \ThesisCite{9}, \ThesisCite{10}, \ThesisCite{11}. شناسایی دقیق این محدودیت‌ها نه تنها مانع از ادعاهای غیرواقع‌بینانه می‌گردد، بلکه نقشه راه شفافی را برای بهینه‌سازی‌های آتی در اختیار سایر پژوهشگران قرار می‌دهد. محدودیت‌های فنی و عملیاتی اصلی این پژوهش به شرح زیر ارزیابی می‌گردند \ThesisCite{1}, \ThesisCite{2}, \ThesisCite{3}, \ThesisCite{9}, \ThesisCite{10}, \ThesisCite{11}:\par

% SOURCE Technical_Boundaries_Chapter5.txt:4-4 [spacing]


% SOURCE Technical_Boundaries_Chapter5.txt:5-5 [paragraph]
۱. \textbf{تمرکز بر مدالیته‌های تصویری دوبعدی ایستا و عدم ارزیابی جریان‌های ویدئویی:}\par

% SOURCE Technical_Boundaries_Chapter5.txt:6-6 [paragraph]
ارزیابی‌های تجربی این پژوهش بر روی ۴ مجموعه‌داده مرجع با مدالیته‌های تصویری دوبعدی ایستا (شامل فاندوس شبکیه \lr{DRIVE/RITE}، اسلایس‌های دوبعدی \lr{BraTS 2020 MRI}، و تصاویر رادیوگرافی \lr{COVID{-}19 CXR}) انجام پذیرفته است \ThesisCite{1}, \ThesisCite{4}, \ThesisCite{9}. بنابر این، عملکرد سیستم بر روی جریان‌های پیوسته ویدئویی (مانند سونوگرافی زمان‌واقعی یا اکوکاردیوگرافی) و حجم‌های سه‌بعدی متراکم \lr{CT/MRI} به صورت مستقیم ارزیابی نشده است \ThesisCite{1}, \ThesisCite{3}.\par

% SOURCE Technical_Boundaries_Chapter5.txt:7-7 [spacing]


% SOURCE Technical_Boundaries_Chapter5.txt:8-8 [paragraph]
۲. \textbf{گلوگاه محاسباتی استنتاج مدل عمیق \lr{U{-}Net} در پردازش لبه:}\par

% SOURCE Technical_Boundaries_Chapter5.txt:9-9 [paragraph]
سنجش زمان اجرای سیستم بر روی پلتفرم کم‌مصرف پردازش لبه \lr{NVIDIA Jetson TX2} نشان داد که زمان اجرای کل برابر با \textbf{۲.۸۵ ثانیه} است \ThesisCite{10}, \ThesisCite{11}. تحلیل تفکیکی زمان اجرا آشکار ساخت که استنتاج شبکه عصبی عمیق \lr{U{-}Net} جهت قطعه‌بندی \lr{ROI} سهمی معادل \textbf{۰.۸۲ ثانیه (۲۸.۸\SourceGlyph{066A} از زمان کل)} را به خود اختصاص داده است \ThesisCite{1}, \ThesisCite{10}. این امر نشان می‌دهد که پیاده‌سازی مدل‌های عمیق بر روی تجهیزات کم‌مصرف پوشیدنی یا حسگرهای محدود \lr{IoMT} می‌تواند با چالش مصرف انرژی و گلوگاه استنتاج مواجه گردد.\par

% SOURCE Technical_Boundaries_Chapter5.txt:10-10 [spacing]


% SOURCE Technical_Boundaries_Chapter5.txt:11-11 [paragraph]
۳. \textbf{محدودیت ظرفیت جای‌دهی فراداده در پس‌زمینه کم‌برجستگی \lr{(Non{-}ROI Capacity Bound)}:}\par

% SOURCE Technical_Boundaries_Chapter5.txt:12-12 [paragraph]
جهت تضمین صیانت ۱۰۰ درصدی از بافت‌های تشخیصی پزشکی ($\mathrm{PSNR}_{\mathrm{ROI}} = \infty$)، ماژول پنهان‌نگاری \lr{Saliency{-}aware} جای‌دهی داده‌ها را صرفاً به ۲ بیت کم‌ارزش \lr{(2{-}bit LSB)} پس‌زمینه کم‌برجستگی محدود نمود \ThesisCite{3}, \ThesisCite{8}, \ThesisCite{9}. اگرچه با به‌کارگیری فشرده‌سازی \lr{zlib} نرخ جای‌دهی به \textbf{۱.۱۸ بیت بر پیکسل \lr{(bpp)}} رسید، اما در تصاویری که ناحیه \lr{ROI} بخش اعظم تصویر را پوشش دهد یا فراداده بیمار بسیار حجیم باشد (مانند فایل‌های صوتی یا ویدئویی سوابق بیمار)، ظرفیت جای‌دهی ممکن است با محدودیت مواجه شود \ThesisCite{3}, \ThesisCite{9}.\par

% SOURCE Technical_Boundaries_Chapter5.txt:13-13 [spacing]


% SOURCE Technical_Boundaries_Chapter5.txt:14-14 [paragraph]
۴. \textbf{الزام وجود کانال امن برای تبادل اولیه کلید عمومی و پارامترهای آشوبی:}\par

% SOURCE Technical_Boundaries_Chapter5.txt:15-15 [paragraph]
تولید کلید هَش ۲۵۶ بیتی $K_{\mathrm{hash}}$ به صورت کاملاً پویا و بر پایه گشتاورهای آماری \lr{ROI} انجام می‌شود؛ اما برای رمزگشایی موفق در سمت گیرنده، پارامترهای کنترل اولیه سیستم ۵بعدی ابرآشوبی یا کلید عمومی باید از طریق یک کانال ارتباطی ایمن (مانند پروتکل تبادل کلید دیفی{-}هلمن یا زیرساخت کلید عمومی \lr{PKI}) میان فرستنده و گیرنده همگام‌سازی شوند \ThesisCite{1}, \ThesisCite{2}, \ThesisCite{4}.\par

% SOURCE Technical_Boundaries_Chapter5.txt:16-16 [spacing]


% SOURCE Technical_Boundaries_Chapter5.txt:17-17 [paragraph]
فرمول‌بندی ریاضی و تبیین ساختاریافته مرزها و محدودیت‌های فنی پژوهش در رابطه ۵{-}۶ ارائه گردیده است \ThesisCite{1}, \ThesisCite{9}, \ThesisCite{10}, \ThesisCite{11}.\par

% SOURCE Technical_Boundaries_Chapter5.txt:18-18 [spacing]


% SOURCE Technical_Boundaries_Chapter5.txt:19-19 [image]
\SourcePicture{eq5_5_technical_boundaries_math.png}{رابطه ۵{-}۶: فرمول‌بندی ریاضی مرزها و محدودیت‌های واقعی فنی پژوهش}{[رابطه ۵{-}۶: فرمول‌بندی ریاضی مرزها و محدودیت‌های واقعی فنی پژوهش {-} \lr{eq5\_5\_technical\_boundaries\_math.png}]}{2}

% SOURCE Technical_Boundaries_Chapter5.txt:20-20 [spacing]

% SOURCE Future_Recommendations_Chapter5.txt:1-1 [heading]
\SourceHeading{1}{۵{-}۷. پیشنهادهای کاربردی و نقشه راه پژوهش‌های آتی \lr{(Practical \& Future Recommendations)}}

% SOURCE Future_Recommendations_Chapter5.txt:2-2 [spacing]


% SOURCE Future_Recommendations_Chapter5.txt:3-3 [paragraph]
بر پایه یافته‌های تجربی حاصل از شبیه‌سازی‌ها، تحلیل‌های آماری و ارزیابی‌های سخت‌افزاری پردازش لبه، پیشنهادهای پژوهشی و نقشه راه توسعه سیستم پیشنهادی در دو بخش مجزای «پیشنهادهای کاربردی جهت استقرار صنعتی و بالینی» و «پیشنهادهای پژوهشی و توسعه علمی آتی» به شرح زیر ارائه می‌گردد \ThesisCite{1}, \ThesisCite{2}, \ThesisCite{9}, \ThesisCite{10}, \ThesisCite{11}:\par

% SOURCE Future_Recommendations_Chapter5.txt:4-4 [spacing]


% SOURCE Future_Recommendations_Chapter5.txt:5-5 [heading]
\SourceHeading{2}{۵{-}۷{-}۱. پیشنهادهای کاربردی جهت استقرار صنعتی و بالینی \lr{(Practical \& Clinical Deployment Recommendations)}:}

% SOURCE Future_Recommendations_Chapter5.txt:6-6 [paragraph]
۱. \textbf{ادغام مستقیم با استانداردهای بین‌المللی پرونده و آرشیو پزشکی \lr{(DICOM / HL7 / HIPAA / PACS)}:}\par

% SOURCE Future_Recommendations_Chapter5.txt:7-7 [paragraph]
جهت عملیاتی‌سازی سیستم در مراکز درمانی و بیمارستان‌های هوشمند، پیاده‌سازی ماژول پنهان‌نگاری \lr{Saliency} و رمزنگاری آشوبی به صورت یک ماژول افزودنی \lr{(Plugin)} مستقیم در استانداردهای آرشیو و تبادل تصاویر پزشکی \lr{(PACS)} تحت رعایت کامل الزامات حفظ حریم خصوصی \lr{HIPAA} و پروتکل‌های ارتباطی \lr{HL7} پیشنهاد می‌گردد \ThesisCite{2}, \ThesisCite{3}, \ThesisCite{9}.\par

% SOURCE Future_Recommendations_Chapter5.txt:8-8 [spacing]


% SOURCE Future_Recommendations_Chapter5.txt:9-9 [paragraph]
۲. \textbf{پیاده‌سازی سخت‌افزاری بر روی تراشه‌های اختصاصی \lr{FPGA} و \lr{ASIC}:}\par

% SOURCE Future_Recommendations_Chapter5.txt:10-10 [paragraph]
جهت دست‌یابی به پردازش سخت‌افزاری موازی زمان‌واقعی در ابعاد میلی‌ثانیه ($T < 10\mathrm{\ ms}$)، پیاده‌سازی الگوریتم ۵بعدی ابرآشوبی و فازهای انتشار \lr{DNA} بر روی تراشه‌های سخت‌افزاری \lr{FPGA} (مانند \lr{Xilinx Zynq UltraScale}+) یا تراشه‌های اختصاصی \lr{ASIC} توصیه می‌شود تا به عنوان کارت‌های شتاب‌دهنده امنیتی درون تجهیزات تصویربرداری (مانند دستگاه‌های سونوگرافی یا \lr{MRI}) تعبیه گردند \ThesisCite{10}, \ThesisCite{11}.\par

% SOURCE Future_Recommendations_Chapter5.txt:11-11 [spacing]


% SOURCE Future_Recommendations_Chapter5.txt:12-12 [heading]
\SourceHeading{2}{۵{-}۷{-}۲. پیشنهادهای پژوهشی و نقشه راه توسعه علمی آتی \lr{(Future Academic \& Research Directions)}:}

% SOURCE Future_Recommendations_Chapter5.txt:13-13 [paragraph]
۱. \textbf{توسعه به جریان‌های ویدئویی پیوسته و مدالیته‌های سه‌بعدی:}\par

% SOURCE Future_Recommendations_Chapter5.txt:14-14 [paragraph]
تعمیم نگاشت ۵بعدی ابرآشوبی و قوانین انتشار \lr{DNA} از تصاویر دوبعدی ایستا به فریم‌های پیوسته ویدیوهای پزشکی (مانند اکوکاردیوگرافی و سونوگرافی زمان‌واقعی) و حجم‌های متراکم ۳بعدی \lr{CT} و \lr{MRI} ($I_{x,\allowbreak{}y,\allowbreak{}z}$) جهت پوشش تمامی مدالیته‌های تصویربرداری بالینی \ThesisCite{1}, \ThesisCite{4}, \ThesisCite{9}.\par

% SOURCE Future_Recommendations_Chapter5.txt:15-15 [spacing]


% SOURCE Future_Recommendations_Chapter5.txt:16-16 [paragraph]
۲. \textbf{تلفیق با رمزشناسی پساکوانتومی \lr{(Post{-}Quantum Cryptography {-} PQC)}:}\par

% SOURCE Future_Recommendations_Chapter5.txt:17-17 [paragraph]
ترکیب نگاشت‌های آشوبی با الگوریتم‌های رمزشناسی مبتنی بر مشبکه‌ (\lr{Lattice{-}based Cryptography} نظیر \lr{CRYSTALS{-}Kyber}) جهت ایمن‌سازی کلید هَش $K_{\mathrm{hash}}$ و همگام‌سازی متغیرهای اولیه در برابر تهدیدات محاسباتی رایانه‌های کوانتومی آتی \ThesisCite{13}.\par

% SOURCE Future_Recommendations_Chapter5.txt:18-18 [spacing]


% SOURCE Future_Recommendations_Chapter5.txt:19-19 [paragraph]
۳. \textbf{بهره‌گیری از معماری‌های سبُک‌تر نظیر \lr{Vision Transformers} \lr{(ViT)} و \lr{MobileNetV4}:}\par

% SOURCE Future_Recommendations_Chapter5.txt:20-20 [paragraph]
جایگزینی یا بهینه‌سازی معماری \lr{U{-}Net} با شبکه‌های سبُک‌تر مبتنی بر مکانیسم توجه \lr{(Attention{-}based)} مانند \lr{Swin Transformer} یا \lr{MobileNetV4} جهت کاهش زمان استنتاج قطعه‌بندی \lr{ROI} به کمتر از ۰.۱۰ ثانیه در حسگرهای پوشیدنی و تجهیزات محدود \lr{IoMT} \ThesisCite{5}, \ThesisCite{10}.\par

% SOURCE Future_Recommendations_Chapter5.txt:21-21 [spacing]


% SOURCE Future_Recommendations_Chapter5.txt:22-22 [paragraph]
فرمول‌بندی ریاضی پیشنهادهای کاربردی و نقشه راه پژوهش‌های آتی در رابطه ۵{-}۷ ارائه گردیده است \ThesisCite{1}, \ThesisCite{10}, \ThesisCite{11}.\par

% SOURCE Future_Recommendations_Chapter5.txt:23-23 [spacing]


% SOURCE Future_Recommendations_Chapter5.txt:24-24 [image]
\SourcePicture{eq5_6_future_roadmap_math.png}{رابطه ۵{-}۷: فرمول‌بندی ریاضی پیشنهادهای کاربردی و نقشه راه پژوهش‌های آتی}{[رابطه ۵{-}۷: فرمول‌بندی ریاضی پیشنهادهای کاربردی و نقشه راه پژوهش‌های آتی {-} \lr{eq5\_6\_future\_roadmap\_math.png}]}{2}

% SOURCE Future_Recommendations_Chapter5.txt:25-25 [spacing]

% SOURCE Final_Remarks_Chapter5.txt:1-1 [heading]
\SourceHeading{1}{۵{-}۸. جمع‌بندی نهایی و پیام کلامی پایان‌نامه \lr{(Final Concluding Remarks)}}

% SOURCE Final_Remarks_Chapter5.txt:2-2 [spacing]


% SOURCE Final_Remarks_Chapter5.txt:3-3 [paragraph]
پژوهش حاضر پاسخ جامع، علمی و پیاده‌سازی‌شده‌ای به چالش حیاتی ایمن‌سازی تصاویر پزشکی و حفظ حریم خصوصی بیماران در بسترهای ارتباطی ناامن شبکه اینترنت اشیای پزشکی \lr{(IoMT)} و سامانه‌های سلامت الکترونیک \lr{(E{-}Health)} ارائه نمود \ThesisCite{1}, \ThesisCite{2}, \ThesisCite{10}. دستاوردهای روش‌شناختی و تجربی این پایان‌نامه اثبات می‌نمایند که بهره‌گیری از پارادایم هیبریدی چندلایه‌ای\SourceGlyph{2014}مبتنی بر تلفیق هوشمند قطعه‌بندی عمیق \lr{U{-}Net}، استخراج کلید پویا \lr{SHA{-}256}، سیستم ۵بعدی ابرآشوبی، جایگشت زیگ‌زاگی و انتشار پویای \lr{DNA}، و پنهان‌نگاری آگاه به برجستگی بصری \lr{(Saliency{-}aware)} همراه با فشرده‌سازی \lr{zlib}\SourceGlyph{2014}می‌تواند تحولی بنیادی در رمزشناسی تصویر پزشکی ایجاد نماید \ThesisCite{1}, \ThesisCite{4}, \ThesisCite{6}, \ThesisCite{8}.\par

% SOURCE Final_Remarks_Chapter5.txt:4-4 [spacing]


% SOURCE Final_Remarks_Chapter5.txt:5-5 [paragraph]
این چارچوب امنیتی ضمن ارائه مقاومت مطلق در برابر تمامی حملات رمزشناسی و سایبری (از جمله حملات تفاضلی با \lr{NPCR=99.6254\%}، حملات متن‌واضح برگزیده \lr{CPA/KPA}، و نویزهای شدید کانال)، اصالت بالینی بافت‌های تشخیصی را به صورت ۱۰۰\SourceGlyph{066A} صیانت نموده ($\mathrm{PSNR}_{\mathrm{ROI}} = \infty$) و توازنی آرمانی میان امنیت، شفافیت دیداری ($\mathrm{PSNR} = 52.41\mathrm{\ dB}$) و کارایی زمان‌واقعی بر روی تجهیزات کم‌مصرف پردازش لبه (زمان اجرای ۲.۸۵ ثانیه بر روی \lr{NVIDIA Jetson TX2}) برقرار می‌سازد \ThesisCite{9}, \ThesisCite{10}, \ThesisCite{11}. پیام کلامی محوری این پایان‌نامه آن است که «امکان ارتقای هم‌زمان امنیت، صیانت بالینی و کارایی عملیاتی بدون قربانی کردن هیچ‌یک از این اضلاع متعارض در سامانه‌های سلامت هوشمند کاملاً میسر و عملیاتی است».\par

% SOURCE Final_Remarks_Chapter5.txt:6-6 [spacing]


% SOURCE Final_Remarks_Chapter5.txt:7-7 [paragraph]
فرمول‌بندی ریاضی پیام کلامی محوری و جمع‌بندی نهایی پایان‌نامه در رابطه ۵{-}۸ نشان داده شده است \ThesisCite{1}, \ThesisCite{10}, \ThesisCite{11}.\par

% SOURCE Final_Remarks_Chapter5.txt:8-8 [spacing]


% SOURCE Final_Remarks_Chapter5.txt:9-9 [image]
\SourcePicture{eq5_7_final_takeaway_math.png}{رابطه ۵{-}۸: فرمول‌بندی ریاضی جمع‌بندی نهایی پایان‌نامه و پارادایم امنیتی \lr{IoMT}}{[رابطه ۵{-}۸: فرمول‌بندی ریاضی جمع‌بندی نهایی پایان‌نامه و پارادایم امنیتی \lr{IoMT {-} eq5\_7\_final\_takeaway\_math.png}]}{2}

% SOURCE Final_Remarks_Chapter5.txt:10-10 [spacing]


% SOURCE Final_Remarks_Chapter5.txt:11-11 [spacing]


% SOURCE Final_Remarks_Chapter5.txt:12-12 [reference-moved]


% SOURCE Final_Remarks_Chapter5.txt:13-13 [reference-moved]


% SOURCE Final_Remarks_Chapter5.txt:14-14 [reference-moved]


% SOURCE Final_Remarks_Chapter5.txt:15-15 [reference-moved]


% SOURCE Final_Remarks_Chapter5.txt:16-16 [reference-moved]


% SOURCE Final_Remarks_Chapter5.txt:17-17 [reference-moved]


% SOURCE Final_Remarks_Chapter5.txt:18-18 [reference-moved]


% SOURCE Final_Remarks_Chapter5.txt:19-19 [reference-moved]


% SOURCE Final_Remarks_Chapter5.txt:20-20 [reference-moved]


% SOURCE Final_Remarks_Chapter5.txt:21-21 [reference-moved]


% SOURCE Final_Remarks_Chapter5.txt:22-22 [reference-moved]


% SOURCE Final_Remarks_Chapter5.txt:23-23 [reference-moved]


% SOURCE Final_Remarks_Chapter5.txt:24-24 [reference-moved]

